\documentclass{beamer}

\usepackage{./styles/esi-beamer}

\title{Chapitre 10: Les tableaux}
\subtitle{\tiny [DEV1] Cours de développement}
\author{\href{https://andrew.sassoye.be}{Andrew SASSOYE}}
\date{2019-2020}

\begin{document}
	\begin{frame}
		\titlepage
	\end{frame}

	\begin{frame}
        \center Document crée par \textbf{Andrew SASSOYE} et produit avec \LaTeX.
        \tiny Version du \today. \normalsize

		\begin{block}{Licence}
            Ce document est mis à disposition sous licence Creative Commons
            Attribution - Partage dans les Mêmes Conditions 4.0 International.

            \begin{center}
                \includegraphics[width=8mm]{./styles/images/cc}
                \includegraphics[width=8mm]{./styles/images/by}
                \includegraphics[width=8mm]{./styles/images/sa}
            \end{center}

            \center\tiny Pour voir une copie de cette licence, visitez
            \href{https://creativecommons.org/licenses/by-sa/4.0/}{https://creativecommons.org/licenses/by-sa/4.0/}
            \normalsize
        \end{block}

        \tiny Le code source de ce document est disponible sur github\\ \href{https://github.com/asassoye/ESI-dev1-slides}{https://github.com/asassoye/ESI-dev1-slides}

    \end{frame}

    \part{Théorie}

    \section{Déclaration et création}\label{sec:declaration-et-creation}
        \subsection{Introduction}\label{subsec:introduction}
    \import{contents/declaration-et-creation/}{/introduction}
        \subsection{Déclaration de tableaux}\label{subsec:declaration-de-tableaux}
    \import{contents/declaration-et-creation/}{declaration-de-tableaux.tex}
        \subsection{Création de tableaux}\label{subsec:creation-de-tableaux}
    \import{contents/declaration-et-creation/}{creation-de-tableaux.tex}

    \section{Utilisation}
        \subsection{Accès individuel aux éléments d'un tableau}\label{subsec:acces-individuel-aux-elements-dun-tableau}
        \subsection{Affectation de tableaux}\label{subsec:affectation-de-tableaux}
        \subsection{La taille d'un tableau: length}\label{subsec:la-taille-dun-tableau}
        \subsection{Boucle for\ldots each}\label{subsec:boucle-foreach}

    \section{Tableau en argument ou en retour}\label{sec:tableau-en-argument-ou-en-retour}

    \section{Tableaux à plusieurs indices}\label{sec:tableaux-a-plusieurs-indices}
        \subsection{Présentation générale}\label{subsec:presentation-generale}
        \subsection{Initialisation}\label{subsec:initialisation}
        \subsection{Exemple}\label{subsec:exemple}
        \subsection{For\ldots each}\label{subsec:for-each}

    \part{Exercices}

\end{document}